% TOP_PROBLEM: {"id": 30007144, "title": "Deforming hypersymplectic triples to hyperkähler triples", "category_id": 6, "category": {"id": 6, "name": "geometry", "display_name": "Geometry", "description": "Euclidean and non-Euclidean geometry, geometric structures.", "slug": "geometry", "order_index": 6, "created_at": "2026-07-31T15:26:25.671Z"}, "status": "open"}
% FIRST_POSTED: 2026-10-06
% ATTEMPT_STATUS: solved
% OPENAI_CLAIM: {"schema_version":1,"source_repo":"https://github.com/openai/math","source_commit":"adc7f1241b42e322a6451854ab7e4b4c146bf78a","release_date":"2026-10-06","match":"full","resolution":"proved","independently_reviewed":false,"second_pass":"agreed","families":[{"family":"341","title":"Donaldson's hypersymplectic deformation conjecture","subject":"Differential geometry","manuscripts":[{"title":"Deforming hypersymplectic four-manifolds to hyperkähler triples","pdf_path":"preprints/Deforming-hypersymplectic-four-manifolds-to-hyperkahler-triples-September-23-2026/paper.pdf","readme_path":"preprints/Deforming-hypersymplectic-four-manifolds-to-hyperkahler-triples-September-23-2026/README.md","date":"2026-09-23","theorem_ref":"cor:individual-kahler (sections/01-introduction.tex:51); thm:main (sections/01-introduction.tex:11)","statement_sources":[{"path":"preprints/Deforming-hypersymplectic-four-manifolds-to-hyperkahler-triples-September-23-2026/build/sections/01-introduction.tex","line":51,"label":"cor:individual-kahler"},{"path":"preprints/Deforming-hypersymplectic-four-manifolds-to-hyperkahler-triples-September-23-2026/build/sections/01-introduction.tex","line":11,"label":"thm:main"}]}],"lean":null,"reasoning_trace":null}]}
\documentclass[11pt]{article}
\usepackage[margin=25mm]{geometry}
\usepackage{fontspec}
\setmainfont{Latin Modern Roman}
\usepackage{amsmath,amssymb,mathtools}
\usepackage{unicode-math}
\setmathfont{Latin Modern Math}
\usepackage{enumitem}
\usepackage{xurl,hyperref,bookmark}
\hypersetup{colorlinks=true,urlcolor=blue}
\setcounter{secnumdepth}{0}
\setlength{\parindent}{0pt}
\setlength{\parskip}{5pt}
\setlength{\emergencystretch}{3em}
\begin{document}
\section*{Deforming hypersymplectic triples to hyperkähler triples}

\subsection{OpenAI's claim}
\textit{Hypersymplectic deformation}: On a closed oriented four-manifold, a positive triple satisfying the paper’s normalization and consisting of closed two-forms is claimed to deform smoothly to a hyperkähler triple while keeping each cohomology class fixed. The starting triple need not already define a hyperkähler metric.

\subsection{Mathematical statement of the claimed result}
\paragraph{Mathematical claim 1.} On an oriented four-manifold, a hypersymplectic structure consists of three real forms $\omega_i$ of degree two, satisfying $d\omega_i=0$, with a pointwise span of dimension three that is positive definite for the wedge pairing specified by the orientation. For a closed connected smooth $X$ and such a triple normalized by $\int_X\omega_i\wedge\omega_j=\delta_{ij}$, the claim gives a smooth hypersymplectic deformation $\omega(t)$ with \[\omega(0)=\omega,\qquad[\omega_i(t)]=[\omega_i],\qquad\omega_i(1)\wedge\omega_j(1)=2\delta_{ij}\mu_1.\] The endpoint forms are parallel self-dual forms for a hyperk\"ahler metric. A further corollary states that each nonzero real combination $\alpha_c=\sum c_i\omega_i$ of the original triple is individually a K\"ahler form for some hyperk\"ahler metric and parallel complex structure, chosen separately for that combination.

\subsection{Outline of the argument}
\textit{Hypersymplectic deformation}: Two forms determine an almost-complex structure, and the third form must keep a taming representative throughout the pair deformation. Positive-current separation and quadratic mass estimates establish a criterion for retaining that representative. An auxiliary fibration and an equal-square pair supply the deformation framework, after which classical Kähler arguments reach the hyperkähler endpoint.

\subsection{Scope relative to this problem}
\textit{Hypersymplectic deformation}: On a closed oriented four-manifold, a positive triple satisfying the paper’s normalization and consisting of closed two-forms is claimed to deform smoothly to a hyperkähler triple while keeping each cohomology class fixed. The starting triple need not already define a hyperkähler metric. The parameters, regularity assumptions and normalization in this question are retained in applying that assertion.

With the hypotheses and conventions above, the claimed result addresses the whole mathematical assertion on this page.

\subsection{Formal verification}
\textit{Hypersymplectic deformation}: The pinned catalogue lists no Lean formalization for this result.

This site has not run the Lean development. These links describe OpenAI’s selected formal statements and should be read with their stated scope.

\subsection{Status caveat}
OpenAI reports that an internal model produced the manuscripts. Their arguments have not been independently verified as proofs by this site. Any portion outside the documented formal scope may contain mistakes, and the catalogue label reports the announced mathematical claim.

\subsection{Sources}
\begin{itemize}
\item OpenAI, \emph{Deforming hypersymplectic four-manifolds to hyperkähler triples}, 2026-09-23; citation key \texttt{OAI:Deforming-hypersymplectic-four-manifolds-to-hyperkahler-triples-September-23-2026}. \href{https://github.com/openai/math/blob/adc7f1241b42e322a6451854ab7e4b4c146bf78a/preprints/Deforming-hypersymplectic-four-manifolds-to-hyperkahler-triples-September-23-2026/paper.pdf}{PDF (GitHub)}; \href{https://github.com/openai/math/raw/adc7f1241b42e322a6451854ab7e4b4c146bf78a/preprints/Deforming-hypersymplectic-four-manifolds-to-hyperkahler-triples-September-23-2026/paper.pdf}{PDF (direct)}; \href{https://github.com/openai/math/blob/adc7f1241b42e322a6451854ab7e4b4c146bf78a/preprints/Deforming-hypersymplectic-four-manifolds-to-hyperkahler-triples-September-23-2026/README.md}{Manuscript page and citation}; \href{https://github.com/openai/math/blob/main/preprints/Deforming-hypersymplectic-four-manifolds-to-hyperkahler-triples-September-23-2026/paper.pdf}{Latest version}; Statement source: \href{https://github.com/openai/math/blob/adc7f1241b42e322a6451854ab7e4b4c146bf78a/preprints/Deforming-hypersymplectic-four-manifolds-to-hyperkahler-triples-September-23-2026/build/sections/01-introduction.tex\#L51}{cor:individual-kahler}; Statement source: \href{https://github.com/openai/math/blob/adc7f1241b42e322a6451854ab7e4b4c146bf78a/preprints/Deforming-hypersymplectic-four-manifolds-to-hyperkahler-triples-September-23-2026/build/sections/01-introduction.tex\#L11}{thm:main}.
\item \href{https://github.com/openai/math/blob/adc7f1241b42e322a6451854ab7e4b4c146bf78a/CONTENTS.md}{OpenAI catalogue entry}.
\end{itemize}
\end{document}
